\section{Cash flow interpretation}

\begin{qfmethod}[Interpreting reported and common-size cash flows]\label{meth:reported-common-size-cash-flows}
\qffield{Statement} First identify the principal sources and uses of cash and whether operations fund investment and distributions. Then reconcile net income to CFO: noncash expenses, working-capital investment, non-operating gains and deferred tax explain why profit differs from cash. Common-size analysis can divide signed net cash flows by revenue, or divide gross inflows and gross outflows by their respective totals. These are different denominators; do not mix them or divide by a near-zero net change in cash \citep[chap.~5, solutions~25--30, PDF pp.~187--188]{statement-analysis-workbook-2020}.
\qffield{Practice} Meridian's CFO 137.5 is 13.75 percent of sales; investing cash flow $-80$ is $-8.0\%$ and financing $-70$ is $-7.0\%$, giving cash change $-1.25\%$. Receipts 970 are the only gross inflow in this simplified year. Operating payments 832.5, capex 80, debt principal 30 and dividends 40 total 982.5 of gross outflows, of which operating payments are 84.73 percent. CFO exceeds net income by 25: depreciation adds 50, while net operating working capital absorbs 25. CFO funds capex, but the remaining 57.5 does not fully fund 70 of debt reduction and dividends; the cash reserve supplies 12.5.
\qffield{Interpretation} A negative investing flow can indicate growth or necessary replacement; positive investing cash from disposals can conceal shrinking capacity. Rising CFO from delaying suppliers or selling receivables is not necessarily sustainable. Before comparing issuers, align interest, dividends, leases and supplier-finance presentation; refer to \hyperref[meth:statement-linked-cash-flows]{the cash-flow classification and conversion entry} for IFRS before/after IFRS 18 and US GAAP. A noncash acquisition or new lease liability does not appear as a cash investing outflow merely because assets increased.
\end{qfmethod}

\begin{qfmethod}[Free cash flow, performance and cash coverage ratios]\label{meth:free-cash-flow-and-cash-coverage}
\qffield{Statement} With interest paid included in CFO, no preferred capital and an immediately usable interest tax shield, $FCFF=CFO+I(1-t)-capex$, and $FCFE=CFO-capex+net\ borrowing$. Equivalently, $FCFF=EBIT(1-t)+D\&A-\Delta NWC-capex$. Use operating, noncash, non-debt working capital consistently. If CFO already excludes interest, add no after-tax interest; instead deduct its tax shield when reversing levered taxes. These are analytical cash-flow definitions, not required statement subtotals \citep[chap.~5, solutions~29--30, PDF p.~188]{statement-analysis-workbook-2020}.
\qffield{Practice} Meridian gives $FCFF=137.5+20(0.75)-80=72.5$, independently $170(0.75)+50-25-80=72.5$. Net borrowing is $-30$, so $FCFE=137.5-80-30=27.5$, also $72.5-15-30$. Dividends of 40 exceed FCFE by 12.5, matching the cash decline. If interest were classified as financing, reported CFO would be 157.5; $157.5-5-80$ still gives FCFF 72.5.

CFO/revenue is 13.75 percent; CFO/net income is 1.222; CFO/capex is 1.719; CFO/dividends is 3.438. A stated pre-interest-and-tax cash interest cover is $(CFO+I+cash\ tax)/I=195/20=9.75$. An after-tax cash debt-service cover is $(CFO+I)/(I+scheduled\ principal)=157.5/50=3.15$. They answer different questions and are not interchangeable with EBIT/interest of 8.5. Coverage above one does not guarantee liquidity on the payment date.
\qffield{Pitfalls} Adding back interest twice, treating net borrowing as operating value creation, omitting replacement capex, or calling FCFE distributable when legal restrictions or required cash reserves prevent a dividend.
\end{qfmethod}
